\documentclass[10pt,a4paper]{amsart}
\usepackage{amsmath,amssymb,amsthm,graphicx}
\usepackage{amsmath}
\usepackage{amssymb}
\usepackage{hyperref}
\usepackage[utf8]{inputenc}
\usepackage{csquotes}
\usepackage[foot]{amsaddr}


\usepackage{times}
\usepackage{amsmath,amsthm,amssymb}
\usepackage{array}
\usepackage{bbm}
\usepackage{xcolor}
\usepackage{verbatim}
\usepackage{url}
%\usepackage{authblk}
\pagestyle{plain}


\def\ve{\varepsilon}
\def\vf{\varphi}
\def\vn{\varnothing}

\def\AA{\mathcal{A}}
\def\AAA{\mathbb{A}}
\def\BB{\mathcal{B}}
\def\CC{\mathcal{C}}
\def\CCC{\mathbb{C}}
\def\DDD{\mathbb{D}}
\def\DD{\mathcal{D}}
\def\FF{\mathcal{F}}
\def\HH{\mathcal{H}}
\def\II{\mathcal{I}}
\def\III{\mathbb{I}}
\def\MM{\mathcal{M}}
\def\NN{\mathcal{N}}
\def\NNN{\mathbb{N}}
\def\PPP{\mathbb{P}}
\def\QQQ{\mathbb{Q}}
\def\RR{\mathcal{R}}
\def\RRR{\mathbb{R}}
\def\SS{\mathcal{S}}
\def\SSS{\mathbb{S}}
\def\UU{\mathcal{U}}
\def\VV{\mathcal{V}}
\def\WW{\mathcal{W}}
\def\ZZ{\mathcal{Z}}
\def\ZZZ{\mathbb{Z}}
\def\PP{\mathcal{P}}
\def\TTT{\mathbb{T}}
\def\EEE{\mathbb{E}}


%%%%  jb abbreviations
\def\bK{{\bf K}}
\def\bL{{\bf L}}
\def\bU{{\bf U}}
\def\aut{\mathop{aut}}
\def\betabar{{\bf\beta}}
\def\Uscr{\mathcal U}
\def\myreduct{\upharpoonright}
\def\abar{{\bf a}}
\def\bbar{\bf b}
\def\cbar{\bf c}
\def\lg{\mathop{\rm lg}}
\def\nmdeg{\mathop{\rm nmdeg}}
\def\icl{\mathop{\rm icl}}
% \sidebar{..} puts .. in comment form
\newcommand{\sidebar}[1]{\vskip10pt\noindent
 \hskip.70truein\vrule width2.0pt\hskip.5em
 \vbox{\hsize= 4truein\noindent\footnotesize\relax #1 }\vskip10pt\noindent}



 %%%%

\def\QQQa{\mathbb{Q}\langle \Bar{\alpha},\alpha_0 \rangle}

\def\cc{\mathfrak{c}}

\def\a{\alpha}
\def\b{\beta}
\def\g{\gamma}
\def\d{\delta}
\def\e{\varepsilon}
\def\z{\zeta}
\def\l{\lambda}
\def\k{\kappa}
\def\o{\theta}
\def\r{\rho}
\def\s{\sigma}
\def\t{\tau}
\def\w{\omega}
\def\D{\Delta}

\def\nmdeg{\operatorname{nmdeg}}
\newcommand{\Fmb}[1]{F_{Mb}\left( #1  \right)}
\newcommand{\indf}[1]{F_{ind}\left( #1  \right)}
\newcommand{\fcb}[1]{F_{cb}\left( #1  \right)}
\newcommand{\ficb}[1]{F_{Icb}\left( #1  \right)}





\DeclareMathOperator{\alg}{alg}
\DeclareMathOperator{\RU}{RU}
\DeclareMathOperator{\RM}{RM}
\DeclareMathOperator{\RD}{RD}
\DeclareMathOperator{\tp}{tp}
\DeclareMathOperator{\acl}{acl}
\DeclareMathOperator{\dcl}{dcl}
\DeclareMathOperator{\Cb}{Cb}
\DeclareMathOperator{\dtr}{\delta-trdeg}


\def\id{\operatorname{id}}
\def\Frac{\operatorname{Frac}}
\def\Const{\operatorname{Const}}
\def\spec{\operatorname{Spec}}
\def\span{\operatorname{span}}
\def\exc{\operatorname{Exc}}
\def\Div{\operatorname{Div}}
\def\cl{\operatorname{cl}}
\def\mer{\operatorname{mer}}
\def\trdeg{\operatorname{trdeg}}
\def\ord{\operatorname{ord}}




\newtheorem{theorem}{Theorem}[section]
\newtheorem{thm}[theorem]{Theorem}
\newtheorem{question}[theorem]{Question}
\newtheorem{lemma}[theorem]{Lemma}
\newtheorem{cor}[theorem]{Corollary}
\newtheorem{prop}[theorem]{Proposition}
\newtheorem{fact}[theorem]{Fact}
\newtheorem{claim}[theorem]{Claim}
\newtheorem{subclaim}[theorem]{Subclaim}
\newtheorem{conj}[theorem]{Conjecture}
\newtheorem{cond}[theorem]{Condition}
\newtheorem{observation}[theorem]{Observation}
\newtheorem{notation}[theorem]{Notation}

\theoremstyle{definition}
\newtheorem{definition}[theorem]{Definition}
\newtheorem{example}[theorem]{Example}
\newtheorem{xca}[theorem]{Exercise}
\newtheorem{ques}[theorem]{Question}
\newtheorem{note}[theorem]{Note}

\theoremstyle{remark}
\newtheorem{remark}[theorem]{Remark}
\newcommand{\m}{\mathbb }
\newcommand{\mc}{\mathcal }
\newcommand{\mf}{\mathfrak }
\newcommand{\mb}{F_{Mb} }
\newcommand{\G}{\mathrm{G} }
\newcommand{\U}{\mathrm{U} }



\def \ta {\tau_{\mathcal{D}/\Delta}}
\def \D {\Delta}
\def \DD {\mathcal D}

\newcommand{\gen}[1]{\left\langle#1\right\rangle}
%\renewcommand{\restriction}{\mathord{\upharpoonright}}

\newcommand{\pd}[2]{\frac{\partial #1}{\partial #2}}
\newcommand{\td}{\text{tr.deg.}}
\newcommand{\pp}{\partial }
\newcommand{\pdtwo}[2]{\frac{\partial^2 #1}{\partial #2^2}}
\newcommand{\od}[2]{\frac{d #1}{d #2}}
\def\Ind{\setbox0=\hbox{$x$}\kern\wd0\hbox to 0pt{\hss$\mid$\hss} \lower.9\ht0\hbox to 0pt{\hss$\smile$\hss}\kern\wd0}
\def\Notind{\setbox0=\hbox{$x$}\kern\wd0\hbox to 0pt{\mathchardef \nn=12854\hss$\nn$\kern1.4\wd0\hss}\hbox to 0pt{\hss$\mid$\hss}\lower.9\ht0 \hbox to 0pt{\hss$\smile$\hss}\kern\wd0}
\def\ind{\mathop{\mathpalette\Ind{}}}
\def\nind{\mathop{\mathpalette\Notind{}}}
\numberwithin{equation}{section}
\newcommand{\is}{^{p^ {-\infty}}}
\newcommand{\QQ}{\mathbb Q}
\newcommand{\fh}{\mathfrak h}
\newcommand{\SL}{\operatorname{SL}}


\title{Simple Homogeneous Structures and Indiscernible Sequence Invariants}
\author[1]{John T. Baldwin}
\author[2]{James Freitag}
\author[3]{Scott Mutchnik}
%\affil[1]{Department of Mathematics, Statistics, and Computer Science, University of Illinois at Chicago}
%%%\affil[2]{Department of Biology, University Y}
%%\author{ John T. Baldwin}
%\author{James Freitag} 
\address{Department of Mathematics, Statistics, and Computer Science, University of Illinois at Chicago }
\email{jbaldwin@uic.edu}
\email{jfreitag@uic.edu}

%\author{Scott Mutchnik}
%\address{Department of Mathematics, Statistics, and Computer Science, University of Illinois at Chicago}
\email{mutchnik@uic.edu}

\thanks{JB was partially supported by Simons travel grant G3535. JF was partially supported by NSF CAREER award number 1945251. SM, who is additionally affiliated with the Institut de mathématiques de Jussieu – Paris Rive Gauche, was supported by the NSF under Grant No. DMS-2303034, and by the European Union’s Horizon 2020 research and innovation programme under the Marie Sklodowska-Curie grant agreement No 101034255.}

\date{\today}
\usepackage{cleveref}
\crefname{originallemma}{original lemma}{original lemmas}
\Crefname{originallemma}{Original lemma}{Original lemmas}
\makeatletter
\newlabel{aclopen}{{4.22}{source}{}{sourcepointer-1}{}}
\newlabel{aclopen@cref}{{[originallemma][0][]4.22}{[1][0][]source}{}{}{}}
\makeatother
\title{Forking independence coincides with tensor independence in the r-space theories of the paper}
\author{John Baldwin, James Freitag, Scott Mutchnik}
\date{arxiv:2405.08211v3; CC BY 4.0}
\begin{document}
\maketitle
\noindent\textit{Editorial note.} All delivered mathematical text is the author's own, copied from the pinned publisher source sources/logic-fl03-simple-homogeneous-structures-and/pinned-original-source.tex; no statement or proof step was rewritten, reordered or removed.\par\medskip
\noindent\textit{Editorial note.} Preamble provenance: the retained preamble is the source's own layout adapter, if any, plus the author's macro block at lines 3-205 of sources/logic-fl03-simple-homogeneous-structures-and/pinned-original-source.tex (everything between the publisher class metadata and \textbackslash{}begin\{document\}, comments excluded). Nothing was added, no macro was invented, and the publisher document class is deliberately not reproduced; the excerpt serves only as preamble, never as a span coordinate.\par\medskip
\noindent\textit{Editorial note.} Metadata provenance: source\_id and source\_version were taken from sources/logic-fl03-simple-homogeneous-structures-and/license-observation.json; the title from SOURCE.json; the author list from the source metadata record; the source URL from the first non-empty of source\_url, source\_urls[0] and the evidence's fixed\_version\_url/license\_url. The licence evidence file used by this card is sources/logic-fl03-simple-homogeneous-structures-and/license-observation.json, recorded in license.evidence\_path so that the binding target is explicit.\par\medskip
\noindent\textit{Editorial note.} External references are navigation pointers only; every deliberately kept step of the author's argument stays inside the counted spans. The displayed numbers of those pointers are the author's own numbering under the author's \textbackslash{}newtheorem declarations (section.index as declared), computed mechanically by the lane's scanner.\par\medskip
\noindent\textit{Editorial note.} Extended pairing, per the Chief's ruling of 2026-10-12: the author's proof environment ends by announcing a claim, so the claim environment (source lines 1124-1128) and the proof that follows it (1133-1169) are counted as the continuation of the same proof and are carried as additional proof spans; the sentence closing the main proof is kept verbatim.\par\medskip
\noindent\textit{Editorial note.} The only reference to an outside item inside the counted spans is Lemma aclopen (proof line 1153), recorded as a navigation stub of kind lemma; the line that closes the proof by referring back to this lemma's own label lies outside the counted ranges.\par\medskip
\noindent\textit{Allowed standard prerequisites (editorial).}\begin{itemize}
\item The two independence notions compared: forking independence over C and the tensor independence defined by free amalgamation of algebraic closures. Set up in the discussion preceding the counted statement, ending at source line 1087. Fixes both sides of the equivalence and the notation used throughout the proof.
\item The r-space theories of the paper and their algebraically closed sets. Used inside the counted proof at lines 1133-1169. Supplies the behaviour of algebraic closures and neighbours that the continuation of the proof uses.
\end{itemize}\medskip
\ifdefined\hypertarget\hypertarget{sourcepointer-1}{}\fi
\noindent\textit{Source reference.} Original lemma 4.22: sources/logic-fl03-simple-homogeneous-structures-and/pinned-original-source.tex, source lines 1153--1153. Navigation pointer for \textbackslash{}ref\{aclopen\} used inside the counted span (line 1153); the argument that consumes it is retained in the span.\par\medskip
\begin{lemma}\label{indeq} $A \ind_{C} B$  if and only if $A
\ind_{C}^{\otimes} B$.
\end{lemma}

\begin{proof}
First, if $A \ind^{\otimes}_{C} B$ with $C \leq M$, then $A \ind_{C} B$
because clearly, $\mathrm{tp}(A/ CB)$ extends to a $C$-invariant global type.
Now suppose that $A \ind_{C} B$. Then, $A \ind_{C} \mathrm{acl}(BC)$, so we
may assume $C \leq B$ and $B$ is algebraically closed. Now, let $A' =
\mathrm{acl}(AC)$ and suppose that $A \nind_{C}^{\otimes} B$. We show that $A
\nind_{C} B$. We may produce a $C$-indiscernible sequence $\{B_{i}\}_{i <
\omega}$ with $B_{0}= B$ so that, for $i < \omega$, $B_{i} \ind_{C}^{\otimes}
\{B_{j}\}_{j < i}$; for $p(X, B) = \mathrm{tp}(A/B)$, it suffices to show
that $\cup \{p(X, B_{i})\}_{i < \omega}$ is inconsistent. Suppose otherwise;
then we may in fact assume that $\{B_{i}\}_{i < \omega}$ is indiscernible
over $A'$. We then get a contradiction from the following claim: \end{proof}

\begin{claim}\label{forcontra} (Here we still assume $A \nind_{C}^{\otimes} B$ with $C \leq B$, $B$
algebraically closed.) Let $\{B_{i}\}_{i < \omega}$ be a
$A'=\mathrm{acl}(AC)$-indiscernible sequence with $B_{0} = B$.
Then $B_{n} \nind_{C}^{\otimes} \{B_{i}\}_{i < n}$.
\end{claim}

\begin{proof} We may assume that $B_{i} \cap B_{j} = C$ for $i \neq j < \omega$;
then clearly, $A' \cap B = C$. It will suffice to show that
$\mathrm{acl}(\{B_{i}\}_{i < n+1}) \supsetneq \cup_{i < n+1} B_{i}$. If $A
\nind_{C}^{\otimes}B$, either $\mathrm{acl}(A'B) = A'B$ and $A'$ and $B$ are
not freely amalgamated over $C$, or  $\mathrm{acl}(A'B) \supsetneq A'B$. In
the first case, there is some $a \in A' \backslash C$ with a neighbor $b$ in
$B \backslash C$. Then $\{a\}$ will have $n+1$ neighbors in $\{B_{i}\}_{i <
n+1}$, so will belong to  $\mathrm{acl}(\{B_{i}\}_{i < n+1})$ but will not
belong to $\cup_{i < n+1} B_{i}$. Therefore, $B_{n} \nind_{C}^{\otimes} \{B_{i}\}_{i < n}$. In the second case, since $\{B_{i}\}_{i <
\omega}$ and thus $\{\mathrm{acl}(A'B_{i})\}_{i < \omega}$ is indiscernible
over $A'$, there is some fixed $D \supseteq A'$ so that
$\mathrm{acl}(A'B_{i}) \cap \mathrm{acl}(A'B_{j}) = D$ for $i \neq j \leq
\omega$, and so that $\{B_{i}\}_{i < \omega}$ is indiscernible over $D$. As
we are in the second case, there are also $\{F_{i}\}_{i < \omega}$, so that
$F_{i} \subseteq \mathrm{acl}(A'B_{i}) \backslash A'B_{i}$ is finite and
closed over $A'B_{i}$ for $i < \omega$, and so that
$\mathrm{acl}(A'B_{i})F_{i} \equiv \mathrm{acl}(A'B_{j})F_{j}$ for $i \neq j
< \omega$. Suppose first that for some (every) $i < \omega$, $F_{i} \subseteq
D$. Each of the $F_{i}$ must have neighbors in $B_{i}$, since $F_i$ is closed
over $A'B_{i}$ and $A'$ is algebraically closed {\bf and so $A'\leq \mathbb{M}$ by
Lemma~\ref{aclopen}}. So if $F_{i} \subseteq D$, then $D$ and $B$ are not
freely amalgamated over $C$, and $B_{n} \nind_{C}^{\otimes} \{B_{i}\}_{i <
n}$ by the first case. Otherwise, for some (every) $i < \omega$, $F_{i}
\not\subseteq D$. There is some finite $D_{0} \subseteq D$ so that, for all
$i < \omega$, $D_{0}$ consists of all neighbors of $F_{i} \backslash D$ in
$D$. Then $F = D_{0} \cup^{n}_{i = 0} (F_{i} \backslash D)$ is a finite set,
not contained in $\cup_{i < n+1} B_{i}$, so that every point of $F$ has at
least $n+1$ neighbors in $F \cup (\cup_{i < n+1} B_{i})$: each point of
$F_{i}$ has $n+1$ neighbors that are either in $D \supseteq A$ (so in $D_{0}
\subseteq F$), in $F_{i} \backslash D$, or in $B_{i}$, while each point of
$D_{0}$ has a neighbor in each of the $F_{i} \backslash D \subseteq F$, and
therefore has at least $n+1$ neighbors in $F$. So $\mathrm{acl}(\{B_{i}\}_{i
< n+1}) \supsetneq \cup_{i < n+1} B_{i}$, because $F \subseteq
\mathrm{acl}(\{B_{i}\}_{i < n+1})$, and again $B_{n} \nind_{C}^{\otimes} \{B_{i}\}_{i < n}$.


\end{proof}

\end{document}
